\section{Results}
\label{sec:results}

\subsection{h-e1: Parameter Validity (PASS)}

\begin{table}[h]
\centering
\caption{Parameter validity results for h-e1 hypothesis.}
\label{tab:he1}
\begin{tabular}{llll}
\toprule
\textbf{Metric} & \textbf{Threshold} & \textbf{Observed} & \textbf{Status} \\
\midrule
NaN/Inf Rate & 0\% & \textbf{0.0\%} & \checkmark PASS \\
Magnitude Ratio & $< 10\times$ & \textbf{0.11$\times$} & \checkmark PASS \\
\bottomrule
\end{tabular}
\end{table}

Duality-derived parameters are numerically stable across all 100 test samples.

\subsection{h-m1: Reconstruction Error (FAIL)}

\begin{table}[h]
\centering
\caption{Reconstruction error comparison for h-m1 hypothesis.}
\label{tab:hm1}
\begin{tabular}{ll}
\toprule
\textbf{Metric} & \textbf{Value} \\
\midrule
Mean Duality Error & 91.45 \\
Mean Random Error & 89.54 \\
Error Reduction & \textbf{-2.13\%} (worse) \\
P-value & $< 0.0001$ \\
Cohen's d & \textbf{-4.56} (large negative) \\
\bottomrule
\end{tabular}
\end{table}

Duality initialization is 2.13\% WORSE than random. The effect is consistent across all 12 BERT layers, with Cohen's d ranging from -4.15 to -5.97.

\subsection{Causal Chain Verification}

\begin{table}[h]
\centering
\caption{Causal chain verification status.}
\label{tab:causal}
\begin{tabular}{lll}
\toprule
\textbf{Step} & \textbf{Description} & \textbf{Status} \\
\midrule
1 & Duality equations $\rightarrow$ valid SSM parameters & \checkmark VERIFIED \\
2 & Duality parameters $\rightarrow$ lower reconstruction error & $\times$ FALSIFIED \\
3 & Lower error $\rightarrow$ better task performance & BLOCKED \\
\bottomrule
\end{tabular}
\end{table}
